\clearpage
\appendix
\section{Vorschriften der Anwendungsfälle im Einzelnen}\label{anh:vorschriften}
\subsection{Fall 1: Berufungsverfahren}\label{anh:fall1}
\offen{Einleitender Satz zu Anhang A gemeinsam formulieren. Die bisherige Referenz der WiMi-Stelle (appendix/anhang\_a\_fall1.tex) ist nicht mehr eingebunden}
% Fall 1 (Berufungsverfahren): Prozess und Referenzzuordnung, Entwurf 19.11.2026, von Hand aus h1/results/entwurf_fall1_berufung übernommen
\paragraph{Prozess} Die Hochschule München besetzt eine Professur in einer Fakultät. Die Hochschulleitung gibt die Professur frei, ein Berufungsausschuss wählt mithilfe eines KI-Tools, von Vorgesprächen, Probevorträgen und Gutachten aus und stellt einen Berufungsvorschlag auf. Die Präsidentin oder der Präsident entscheidet, erteilt den Ruf und lässt die ausgewählte Person nach der Berufungsverhandlung ernennen. Die übrigen Bewerbenden erhalten eine Absage, und ihre Daten werden nach einer festgelegten Frist gelöscht.
\begin{small}
\begin{longtable}{@{}>{\raggedright\arraybackslash}p{1.0cm}>{\raggedright\arraybackslash}p{3.6cm}>{\raggedright\arraybackslash}p{10.0cm}@{}}
\caption[Prozess für Fall~1: Berufungsverfahren]{Prozess für Fall~1: Berufungsverfahren mit Teilprozessen, Tasks und Beschreibungen. Quelle: eigene Darstellung (Entwurf).}\label{tab:fall1_prozess}\\
\toprule \textbf{Kennung} & \textbf{Bezeichnung} & \textbf{Beschreibung} \\ \midrule \endfirsthead
\toprule \textbf{Kennung} & \textbf{Bezeichnung} & \textbf{Beschreibung} \\ \midrule \endhead
\bottomrule \endfoot
\multicolumn{3}{@{}l}{\textbf{TP1 Professur freigeben und ausschreiben}} \\
T1.1 & Bedarf anmelden & Die Fakultät beantragt bei der Hochschulleitung die Wiederbesetzung einer frei werdenden Professur und schlägt die fachliche Ausrichtung vor. \\
T1.2 & Professur freigeben & Die Hochschulleitung entscheidet nach Anhörung des Fakultätsrats über Besetzung und Ausrichtung der Professur. \\
T1.3 & Berufungsausschuss bilden & Der Fakultätsrat setzt einen Ausschuss aus Professorinnen und Professoren, je einer Vertretung der Mitarbeitenden und der Studierenden, der Gleichstellungsbeauftragten und einem auswärtigen Mitglied ein. \\
T1.4 & Ausschreibung veröffentlichen & Die Personalabteilung veröffentlicht die Ausschreibung auf dem Stellenportal der Hochschule und in internationalen Fachportalen. \\
\multicolumn{3}{@{}l}{\textbf{TP2 Bewerbungen annehmen}} \\
T2.1 & Unterlagen empfangen & Bewerbende laden Anschreiben, Lebenslauf, Publikationsliste und Lehrkonzept im Bewerbungsportal hoch. \\
T2.2 & Eingang bestätigen & Das Portal sendet jeder bewerbenden Person eine Eingangsbestätigung per E-Mail. \\
T2.3 & Bewerbungen anlegen & Die Personalabteilung legt für jede Bewerbung einen Datensatz im Bewerbermanagementsystem an. \\
\multicolumn{3}{@{}l}{\textbf{TP3 KI-Vorauswahl}} \\
T3.1 & Unterlagen auswerten & Das KI-Tool liest die Unterlagen und gleicht sie mit den Anforderungen der Ausschreibung ab. \\
T3.2 & Ranking erzeugen & Das KI-Tool ordnet die Bewerbungen nach dem Grad der Übereinstimmung in einer Rangliste. \\
T3.3 & Ranking prüfen & Der Berufungsausschuss prüft Rangliste und Unterlagen und passt die Reihenfolge bei Bedarf an. \\
T3.4 & Shortlist festlegen & Der Ausschuss wählt die Bewerbenden für die Vorgespräche aus. \\
\multicolumn{3}{@{}l}{\textbf{TP4 Vorgespräche und Probevorträge}} \\
T4.1 & Zu Vorgesprächen einladen & Die Personalabteilung lädt die Bewerbenden der Shortlist per E-Mail zu einem persönlichen Vorgespräch ein. \\
T4.2 & Vorgespräche führen & Mindestens zwei Mitglieder des Ausschusses führen mit jeder eingeladenen Person ein Gespräch von 30 bis 60 Minuten. \\
T4.3 & Zum Probevortrag einladen & Die Personalabteilung lädt die nach den Vorgesprächen ausgewählten Bewerbenden per E-Mail zu Probevortrag und Gespräch ein. \\
T4.4 & Probevortrag und Gespräch durchführen & Die Kandidierenden halten eine Lehrprobe vor Ausschuss und Studierenden und führen anschließend ein Gespräch mit dem Ausschuss. \\
T4.5 & Ergebnisse dokumentieren & Der Ausschuss hält Verlauf und Bewertung jedes Vortrags und Gesprächs schriftlich fest. \\
\multicolumn{3}{@{}l}{\textbf{TP5 Berufungsvorschlag}} \\
T5.1 & Gutachten einholen & Der Ausschuss beauftragt auswärtige Professorinnen oder Professoren mit vergleichenden Gutachten zu den Kandidierenden der engeren Wahl. \\
T5.2 & Berufungsliste beschließen & Der Ausschuss beschließt eine begründete Liste mit drei Namen in einer Reihenfolge. \\
T5.3 & Über die Berufung entscheiden & Die Präsidentin oder der Präsident entscheidet anhand der Liste, wer den Ruf erhält. \\
\multicolumn{3}{@{}l}{\textbf{TP6 Ruf und Ernennung}} \\
T6.1 & Ruf erteilen & Die Hochschule teilt der ausgewählten Person schriftlich mit, dass sie berufen werden soll. \\
T6.2 & Berufungsverhandlung führen & Hochschulleitung und berufene Person vereinbaren Ausstattung, Dienstbeginn und Leistungsbezüge. \\
T6.3 & Ernennen & Die Hochschule ernennt die Person durch Aushändigung einer Urkunde in ein Beamtenverhältnis. \\
\multicolumn{3}{@{}l}{\textbf{TP7 Absagen und Löschen}} \\
T7.1 & Absagen versenden & Die Personalabteilung sendet den übrigen Bewerbenden eine Absage per E-Mail. \\
T7.2 & Aufbewahrungsfrist abwarten & Die Personalabteilung bewahrt die Unterlagen der Abgelehnten bis zum Ende einer festgelegten Frist auf. \\
T7.3 & Daten löschen & Die Personalabteilung löscht die Daten im Bewerbungsportal und im Bewerbermanagementsystem. \\
\end{longtable}
\end{small}
\begin{small}
\begin{longtable}{@{}>{\raggedright\arraybackslash}p{1.6cm}>{\raggedright\arraybackslash}p{2.6cm}>{\raggedright\arraybackslash}p{5.6cm}>{\raggedright\arraybackslash}p{2.4cm}>{\raggedright\arraybackslash}p{2.0cm}@{}}
\caption[Referenzzuordnung für Fall~1: Berufungsverfahren]{Referenzzuordnung für Fall~1: Berufungsverfahren. Angaben mit dem Vermerk \enquote{prüfen} sind noch am amtlichen Text zu kontrollieren. Quelle: eigene Darstellung (Entwurf).}\label{tab:fall1_referenz}\\
\toprule \textbf{Regelwerk} & \textbf{Fundstelle} & \textbf{Inhalt} & \textbf{Tasks} & \textbf{Art} \\ \midrule \endfirsthead
\toprule \textbf{Regelwerk} & \textbf{Fundstelle} & \textbf{Inhalt} & \textbf{Tasks} & \textbf{Art} \\ \midrule \endhead
\bottomrule \endfoot
GG & Art. 33 Abs. 2 & Bestenauslese im öffentlichen Dienst & gesamter Prozess & Zuordnung \\
AGG & §~7 & Benachteiligungsverbot, für Beamte über §~24 AGG & gesamter Prozess & Zuordnung \\
DSGVO & Art. 32 & Sicherheit der Verarbeitung (Portal, System, Aufbewahrung) & gesamter Prozess & Zuordnung \\
BayHIG & Art. 66 Abs. 1 & Freigabe und Ausrichtung durch die Hochschulleitung, Anhörung Fakultätsrat & T1.1, T1.2 & Zuordnung \\
BayHIG & Art. 66 Abs. 4 & Bildung und Zusammensetzung des Berufungsausschusses & T1.3 & Zuordnung \\
BayHIG & Art. 66 Abs. 4 Satz 1 Nr. 3 & Gleichstellungsbeauftragte im Ausschuss (Kontrollschritt) & T1.3 & Zuordnung \\
BayHIG & Art. 66 Abs. 3 & öffentliche, in der Regel internationale Ausschreibung & T1.4 & Zuordnung \\
AGG & §~11 & diskriminierungsfreie Ausschreibung & T1.4 & Zuordnung \\
BayHIG & Art. 57 Abs. 3 & Einstellungsvoraussetzungen an einer HAW: mindestens fünf Jahre Berufspraxis, davon drei außerhalb der Hochschule & T1.4, T3.1, T5.2 & Zuordnung \\
DSGVO & Art. 6 & Rechtsgrundlage der Verarbeitung von Bewerberdaten & T2.1, T2.3 & Zuordnung \\
BayDSG & (Norm prüfen) & Verarbeitung nach bayerischem Landesrecht & T2.1, T2.3 & Zuordnung \\
DSGVO & Art. 13 & Informationspflicht bei Erhebung & T2.1, T2.2 & Zuordnung \\
AI Act & Anhang III Nr. 4 & Einstufung als Hochrisiko-KI-System & T3.1, T3.2 & Zuordnung \\
AI Act & Art. 26 & Betreiberpflichten: Verwendung gemäß Anleitung, menschliche Aufsicht, Protokollierung & T3.1, T3.2, T3.3, T3.4 & Zuordnung \\
AGG & §~7 & algorithmische Benachteiligung & T3.1, T3.2 & Zuordnung \\
DSGVO & Art. 22 & keine ausschließlich automatisierte Einzelentscheidung & T3.4, T5.3 & Zuordnung \\
SGB IX & §~165 & Einladung schwerbehinderter Bewerbender (öffentlicher Arbeitgeber) & T4.1, T4.3 & Zuordnung \\
AGG & §~7 & zulässige Fragen im Gespräch & T4.2, T4.4 & Zuordnung \\
DSGVO & Art. 5 Abs. 2 & Rechenschaftspflicht Dokumentation & T4.5, T7.2 & Zuordnung \\
BayHIG & Art. 66 Abs. 5 & auswärtige und vergleichende Gutachten, Vorschlag mit drei Namen, Begründung der Reihung & T5.1, T5.2 & Zuordnung \\
BayHIG & Art. 66 Abs. 6 & Entscheidung der Präsidentin oder des Präsidenten ohne Bindung an die Reihung & T5.3 & Zuordnung \\
BayHIG & Art. 66 Abs. 11 & Ausstattung grundsätzlich befristet & T6.2 & Zuordnung \\
BayBesG & Art. 70 Abs. 1 & Berufungs-Leistungsbezüge aus Anlass der Berufungsverhandlung & T6.2 & Zuordnung \\
BeamtStG & §~8 & Ernennung durch Aushändigung einer Urkunde & T6.3 & Zuordnung \\
BayHIG & Art. 58 Abs. 1 & in der Regel Beamtenverhältnis auf Lebenszeit, ohne drei Jahre Hochschultätigkeit zunächst auf Probe & T6.3 & Zuordnung \\
AGG & §~15 Abs. 4 & 2-Monats-Frist, daraus Aufbewahrung & T7.2 & Zuordnung \\
DSGVO & Art. 17 & Löschung & T7.3 & Zuordnung \\
BayHIG & Art. 66 Abs. 1 Satz 6 & Staatsministerium vor der Ausschreibung informieren & -- & fehlender Schritt \\
BayHIG & Art. 66 Abs. 5 Satz 8 & Stellungnahme des Senats zum Berufungsvorschlag & -- & fehlender Schritt \\
SGB IX & §~164 & Prüfung Besetzung mit schwerbehinderten Menschen & -- & fehlender Schritt \\
SGB IX & §~165 & Meldung der freien Stelle an die Agentur für Arbeit & -- & fehlender Schritt \\
SGB IX & §~178 & Beteiligung der Schwerbehindertenvertretung & -- & fehlender Schritt \\
AI Act & Art. 26 Abs. 11 & Information der Betroffenen über den KI-Einsatz & -- & fehlender Schritt \\
AI Act & Art. 27 & Grundrechte-Folgenabschätzung (öffentliche Stelle) & -- & fehlender Schritt \\
AI Act & Art. 4 & KI-Kompetenz der Mitarbeitenden & -- & fehlender Schritt \\
DSGVO & Art. 35 & Datenschutz-Folgenabschätzung & -- & fehlender Schritt \\
BayPVG & Art. 75a Abs. 1 Nr. 2 & Mitbestimmung des Personalrats bei Einführung und Anwendung des KI-Tools als automatisiertes Verfahren zur Personalverwaltung (Anwendbarkeit auf Berufungen prüfen) & -- & fehlender Schritt \\
DSGVO & Art. 5 Abs. 1 & Grundsätze der Verarbeitung & gesamter Prozess & neutral \\
DSGVO & Art. 9 & besondere Kategorien personenbezogener Daten & gesamter Prozess & neutral \\
DSGVO & Art. 12 & Form der Information & gesamter Prozess & neutral \\
DSGVO & Art. 24 & Verantwortung des Verantwortlichen & gesamter Prozess & neutral \\
DSGVO & Art. 25 & Datenschutz durch Technikgestaltung & gesamter Prozess & neutral \\
DSGVO & Art. 28 & Auftragsverarbeitung (KI-Anbieter) & gesamter Prozess & neutral \\
DSGVO & Art. 30 & Verzeichnis von Verarbeitungstätigkeiten & gesamter Prozess & neutral \\
AI Act & Art. 14 & menschliche Aufsicht (Anbieterpflicht, über Art. 26 relevant) & gesamter Prozess & neutral \\
AI Act & Art. 49 & Registrierung (öffentliche Betreiber) & gesamter Prozess & neutral \\
AI Act & Art. 86 & Recht auf Erläuterung & gesamter Prozess & neutral \\
BayHIG & Art. 66 Abs. 2 & Berichterstatterin oder Berichterstatter ("in der Regel") & gesamter Prozess & neutral \\
BayHIG & Art. 22 & Gleichstellungsbeauftragte, Aufgaben & gesamter Prozess & neutral \\
BayGlG & -- & allgemeines Gleichstellungsrecht, an Hochschulen durch BayHIG überlagert (prüfen) & gesamter Prozess & neutral \\
HM-Satzung & Berufungsordnung der HM (bisher nicht gefunden) & Ebene unterhalb des Gesetzes & gesamter Prozess & neutral \\
BayPVG & Art. 4 Abs. 4 Buchst. a & Professorinnen und Professoren nach Art. 58 BayHIG gelten nicht als Beschäftigte, daher keine Mitbestimmung bei ihrer Einstellung. Meldung als Lücke ist eine Fehlmeldung & -- & Falle \\
\end{longtable}
\end{small}

\subsection{Fall 2: Brandmeldetechnik}
\offen{Vorschriften Fall 2, sobald die Referenzzuordnung erstellt ist}
\subsection{Fall 3: TETRA-Endgeräte}
\offen{Vorschriften Fall~3, sobald die Referenzzuordnung erstellt ist}
\subsection{Fall 4: Vergleichsdaten von Sai et al.}
\offen{Verweis auf die Daten der Studie (h1/input/sai) und Kurzübersicht der Dokumente}

\section{Prüfbogen für die fachliche Prüfung}\label{anh:pruefbogen}
\offen{Entwurf 20.11.2026, Formulierungen gemeinsam durchgehen. Der Bogen wird je Fall mit der Tabelle der Referenz verschickt}
Vielen Dank, dass Sie die Referenzzuordnung prüfen.
Die Referenzzuordnung legt fest, welche Vorschriften ein fachkundiger Mensch dem beschriebenen Prozess zuordnen würde.
Ein KI-Agent wird später an dieser Zuordnung gemessen.
Die Prüfung dauert etwa 45 bis 60 Minuten.

\subsection*{Teil 1: Prüfung der einzelnen Zeilen}
Bitte prüfen Sie jede Zeile der Referenzzuordnung mit drei Fragen.
\begin{enumerate}
  \item Gilt die Vorschrift für diesen Prozess?
  \item Gehört sie zu den genannten Tasks?
  \item Stimmt die Art der Zeile (Zuordnung, fehlender Schritt, neutral, Falle)?
\end{enumerate}
Tragen Sie bitte nur Zeilen ein, bei denen Sie mindestens eine Frage mit \enquote{nein} oder \enquote{unsicher} beantworten.

\begin{small}
\begin{longtable}{@{}>{\raggedright\arraybackslash}p{1.6cm}>{\raggedright\arraybackslash}p{1.6cm}>{\raggedright\arraybackslash}p{1.8cm}>{\raggedright\arraybackslash}p{1.8cm}>{\raggedright\arraybackslash}p{7.4cm}@{}}
\toprule \textbf{Zeile} & \textbf{Gilt?} & \textbf{Task richtig?} & \textbf{Art richtig?} & \textbf{Korrektur und Begründung} \\ \midrule \endhead
\bottomrule \endfoot
 & & & & \\[1.2em]
 & & & & \\[1.2em]
 & & & & \\[1.2em]
 & & & & \\[1.2em]
 & & & & \\[1.2em]
\end{longtable}
\end{small}

\subsection*{Teil 2: Allgemeine Fragen}
\begin{enumerate}
  \item Welche Vorschriften fehlen in der Referenz? Bitte mit Regelwerk, Fundstelle und Task.
  \item Sind die als fehlend markierten Schritte tatsächlich vorgeschrieben?
  \item Entspricht der Prozess dem üblichen Ablauf in der Praxis, abgesehen von den bewusst fehlenden Schritten?
  \item Sind die Annahmen zum Fall realistisch?
  \item Stimmen Kontrollschritt und Falle?
\end{enumerate}

\subsection*{Teil 3: Angaben zur prüfenden Person}
Rolle, Fachgebiet, Jahre Berufserfahrung im Fachgebiet, Datum der Prüfung.
Die Arbeit nennt Sie ausschließlich mit Rolle und Fachgebiet, sofern Sie keiner Namensnennung zustimmen.

\section{Testaufbau der Hypothesen}\label{anh:testaufbau}
% Sätze 1-3 abgesegnet am 11.11.2026
Dieser Anhang zeigt für jede Hypothese, welche Eingaben der Agent erhält und womit sein Ergebnis verglichen wird.
Die Angaben gelten für den Stand der Versuche, auf dem die Ergebnisse in Kapitel~6 beruhen.
\offen{Satz: Alle genannten Dateien liegen im öffentlichen Repository \url{https://github.com/SebastianHaug35/Masterthesis_Experiment}}

\subsection{H1: Regelwerke bestimmen (Fall 1)}\label{anh:testaufbau_h1}
Für H1 erhält der Agent zwei Eingaben, das Profil der Organisation und den Prozess.
% automatisch erzeugt von tools/anhang_b.py -- nicht von Hand ändern

\subsubsection{Eingabe}
\offen{Satz: Verweis auf Tabelle~\ref{tab:h1_profil} (Organisationsprofil)}
\begin{small}
\begin{longtable}{@{}>{\raggedright\arraybackslash}p{5.0cm}>{\raggedright\arraybackslash}p{9.5cm}@{}}
\caption[Organisationsprofil für Fall~1]{Organisationsprofil für Fall~1. Quelle: eigene Darstellung, Daten aus \repodatei{h1/input/firmenprofil.json}.}\label{tab:h1_profil}\\
\toprule \textbf{Merkmal} & \textbf{Angabe} \\ \midrule \endfirsthead
\toprule \textbf{Merkmal} & \textbf{Angabe} \\ \midrule \endhead
\bottomrule \endfoot
Name & Hochschule München \\
Land & Deutschland \\
Bundesland & Bayern \\
Rechtsform & Körperschaft des öffentlichen Rechts / staatliche Einrichtung (bayerische Hochschule) \\
Sektor & öffentlich \\
Branche & Bildung und Forschung \\
Beschäftigte & über 1.000 \\
Personalvertretung & Personalrat \\
Gleichstellungsbeauftragte & ja \\
KI-Einsatz & KI-Tool zur Vorauswahl von Bewerbungen (Ranking) \\
Dienstleister außerhalb der EU & nein \\
Stelle & Wissenschaftliche Mitarbeiterin/wissenschaftlicher Mitarbeiter nach TV-L \\
Befristung & befristet auf 3 Jahre, Qualifizierungsstelle (Promotion) \\
\end{longtable}
\end{small}

\offen{Satz: Verweis auf Tabelle~\ref{tab:h1_prozess} (Prozess mit 6 Teilprozessen und 19 Tasks)}
\begin{small}
\begin{longtable}{@{}>{\raggedright\arraybackslash}p{4.2cm}>{\raggedright\arraybackslash}p{1.3cm}>{\raggedright\arraybackslash}p{9.0cm}@{}}
\caption[Prozess für Fall~1: Bewerbungsprozess wissenschaftliche Mitarbeiterstelle]{Prozess für Fall~1: Bewerbungsprozess wissenschaftliche Mitarbeiterstelle. Quelle: eigene Darstellung, Daten aus \repodatei{h1/input/prozess.json}.}\label{tab:h1_prozess}\\
\toprule \textbf{Teilprozess} & \textbf{Task} & \textbf{Bezeichnung} \\ \midrule \endfirsthead
\toprule \textbf{Teilprozess} & \textbf{Task} & \textbf{Bezeichnung} \\ \midrule \endhead
\bottomrule \endfoot
TP1 Stelle ausschreiben & T1.1 & Bedarf melden \\
 & T1.2 & Stellenprofil und Eingruppierung nach TV-L festlegen \\
 & T1.3 & Ausschreibung intern und extern veröffentlichen \\
TP2 Bewerbungen annehmen & T2.1 & Unterlagen über Bewerbungsportal empfangen \\
 & T2.2 & Eingang bestätigen \\
 & T2.3 & Bewerberdaten im System anlegen \\
TP3 KI-Vorauswahl & T3.1 & Unterlagen mit KI-Tool auswerten \\
 & T3.2 & Ranking erzeugen \\
 & T3.3 & Ergebnis durch Fachabteilung prüfen \\
 & T3.4 & Shortlist festlegen \\
TP4 Vorstellungsgespräche & T4.1 & Kandidaten einladen (inkl. Hinweis auf das Antragsrecht zur Beteiligung der Gleichstellungsbeauftragten) \\
 & T4.2 & Gespräche mit der Auswahlkommission führen (Gleichstellungsbeauftragte auf Antrag der Bewerbenden beteiligt) \\
 & T4.3 & Gespräche dokumentieren \\
TP5 Auswahl und Einstellung & T5.1 & Auswahlentscheidung treffen \\
 & T5.2 & Vertragsangebot erstellen \\
 & T5.3 & Arbeitsvertrag abschließen \\
TP6 Absagen und Löschen & T6.1 & Absagen versenden \\
 & T6.2 & Aufbewahrungsfrist abwarten \\
 & T6.3 & Bewerberdaten löschen \\
\end{longtable}
\end{small}

\offen{Satz: Auftrag an den Agenten (Prompt Stufe~1), Verweis auf \repodatei{h1/prompts/stufe1.md}. Entscheiden: Prompt vollständig abdrucken oder nur verlinken}

\subsubsection{Referenz}
\offen{Satz: Herleitung der Referenz aus der Referenzzuordnung (Anhang~\ref{anh:vorschriften}): es zählen Regelwerke mit mindestens einer Zeile der Art Zuordnung oder fehlender Schritt}
\begin{small}
\begin{longtable}{@{}>{\raggedright\arraybackslash}p{3.0cm}>{\raggedright\arraybackslash}p{3.0cm}>{\raggedleft\arraybackslash}p{3.2cm}>{\raggedleft\arraybackslash}p{3.8cm}@{}}
\caption[Referenz für H1 in Fall~1: 12 Regelwerke]{Referenz für H1 in Fall~1: 12 Regelwerke. Quelle: eigene Darstellung, Daten aus \repodatei{h1/input/soll_liste.csv}.}\label{tab:h1_referenz}\\
\toprule \textbf{Ebene} & \textbf{Regelwerk} & \textbf{Zeilen Zuordnung} & \textbf{Zeilen fehlender Schritt} \\ \midrule \endfirsthead
\toprule \textbf{Ebene} & \textbf{Regelwerk} & \textbf{Zeilen Zuordnung} & \textbf{Zeilen fehlender Schritt} \\ \midrule \endhead
\bottomrule \endfoot
EU & AI Act & 2 & 3 \\
EU & DSGVO & 6 & 1 \\
Bund & AGG & 5 & 0 \\
Bund & GG & 1 & 0 \\
Bund & NachwG & 1 & 0 \\
Bund & SGB IX & 1 & 3 \\
Bund & TzBfG & 1 & 0 \\
Bund & WissZeitVG & 1 & 0 \\
Land Bayern & BayDSG & 1 & 0 \\
Land Bayern & BayGlG & 2 & 0 \\
Land Bayern & BayPVG & 0 & 2 \\
Tarifvertrag & TV-L & 2 & 0 \\
\end{longtable}
\end{small}

\offen{Satz: neutrale Regelwerke (BayHIG) zählen weder als Treffer noch als Fehler; weitere vom Agenten genannte Regelwerke zählen als zusätzlich und senken die Genauigkeit}
\offen{Tabelle zeigt noch die WiMi-Stelle (12 Regelwerke). Neu erzeugen, sobald die Eingabedateien per Auftrag aktualisiert sind: Fall~1 Berufungsverfahren mit 10 Regelwerken, dazu Tabellen für Fall~2 (9) und Fall~3 (7). Schwelle ist 0,84 (Abschnitt~\ref{sec:vergleichswerte}); Skript tools/anhang\_b.py entsprechend anpassen}

\offen{Fall~2 bis Fall~4 für H1 ergänzen, sobald die Eingaben vorliegen}

\subsection{H2: Vorschriften zuordnen}
\offen{Eingabe wie H1 plus Rechtsrahmen (Konfiguration A: aus Stufe~1, Konfiguration B: vorgegeben, \repodatei{h1/input/rechtsrahmen_b.json}). Referenz: Zeilen der Art Zuordnung in Anhang~\ref{anh:vorschriften}}

\subsection{H3: Fehlende Schritte erkennen}
\offen{Eingabe wie H2. Referenz: Zeilen der Art fehlender Schritt in Anhang~\ref{anh:vorschriften}, Kontrolle Gleichstellungsbeauftragte}

\subsection{H4: Prozess selbst entwerfen}
\offen{Einleitende Sätze zu H4 (Eingabe ohne Prozess, Bewertung über Reihenfolge-Paare) gemeinsam formulieren}
\begin{small}
\begin{longtable}{@{}>{\raggedright\arraybackslash}p{0.7cm}>{\raggedright\arraybackslash}p{4.8cm}>{\raggedright\arraybackslash}p{4.4cm}>{\raggedright\arraybackslash}p{4.5cm}@{}}
\caption[Reihenfolge-Paare für Fall~1]{Reihenfolge-Paare \enquote{A vor B} für Fall~1 (Berufungsverfahren). Schritte ohne Task-Kennung fehlen im vorgegebenen Prozess. Quelle: eigene Darstellung (Entwurf).}\label{tab:h4_paare_fall1}\\
\toprule \textbf{Nr.} & \textbf{Schritt A} & \textbf{Schritt B} & \textbf{Vorschrift} \\ \midrule \endfirsthead
\toprule \textbf{Nr.} & \textbf{Schritt A} & \textbf{Schritt B} & \textbf{Vorschrift} \\ \midrule \endhead
\bottomrule \endfoot
1 & Professur freigeben (T1.2) & Ausschreibung veröffentlichen (T1.4) & Art.~66 Abs.~1 BayHIG \\
2 & Staatsministerium informieren & Ausschreibung veröffentlichen (T1.4) & Art.~66 Abs.~1 Satz~6 BayHIG \\
3 & Besetzung mit schwerbehinderten Menschen prüfen & Ausschreibung veröffentlichen (T1.4) & §~164 SGB~IX \\
4 & Stelle der Agentur für Arbeit melden & Shortlist festlegen (T3.4) & §~165 SGB~IX \\
5 & Berufungsausschuss mit Gleichstellungsbeauftragter bilden (T1.3) & Ranking prüfen (T3.3) & Art.~66 Abs.~4 BayHIG \\
6 & Bewerbende über den KI-Einsatz informieren & Unterlagen auswerten (T3.1) & Art.~26 Abs.~11 AI Act, Art.~13 DSGVO \\
7 & Datenschutz-Folgenabschätzung durchführen & Unterlagen auswerten (T3.1) & Art.~35 DSGVO \\
8 & Grundrechte-Folgenabschätzung durchführen & Unterlagen auswerten (T3.1) & Art.~27 AI Act \\
9 & Mitarbeitende im Umgang mit KI schulen & Unterlagen auswerten (T3.1) & Art.~4 AI Act \\
10 & Ranking prüfen (T3.3) & Shortlist festlegen (T3.4) & Art.~26 AI Act, Art.~22 DSGVO \\
11 & Schwerbehindertenvertretung beteiligen & Shortlist festlegen (T3.4) & §~178 Abs.~2 SGB~IX \\
12 & Gutachten einholen (T5.1) & Berufungsliste beschließen (T5.2) & Art.~66 Abs.~5 BayHIG \\
13 & Stellungnahme des Senats einholen & Über die Berufung entscheiden (T5.3) & Art.~66 Abs.~5 Satz~8 BayHIG \\
14 & Ruf erteilen (T6.1) & Ernennen (T6.3) & Art.~66 Abs.~6 BayHIG, §~8 BeamtStG \\
15 & Absagen versenden (T7.1), danach Frist abwarten & Daten löschen (T7.3) & §~15 Abs.~4 AGG, Art.~17 DSGVO \\
\end{longtable}
\end{small}
\offen{Mögliches Paar 16: Personalrat beteiligen vor Einführung des KI-Werkzeugs (Art.~75a Abs.~1 Nr.~2 BayPVG). Anwendbarkeit auf Berufungen durch Prof.~Böhm prüfen lassen, liegt nah an der Falle}
\offen{Reihenfolge-Paare für Fall~2 und Fall~3}

\subsection{H5 bis H7: Prüfregeln und Laufzeit-Kontrolle}
\offen{Referenzfälle mit Soll-Entscheidung, noch nicht erstellt}
